% ============================================================
% f2014_beijing_arts_q11_threeviews.tex
% 2014 北京卷（文科）第 11 题：三棱锥的三视图
% 来源 PDF：普通高考/2014/2014北京文.pdf
% 原页：第 1 页（A3），印面 863，第 11 题题干下方
% 正（主）视图：底 2、高 2 的直角三角形，另有自顶点引向底边中点的投影线；
% 侧（左）视图：底 1、高 2 的直角三角形；俯视图：底 2、顶点在底边中点的三角形。
% ============================================================
\begin{tikzpicture}[>=Stealth, line width=0.6pt, font=\small]
  % ---------- 正（主）视图 ----------
  \draw (0,0) -- (3,0) -- (0,3) -- cycle;
  \draw (0,3) -- (1.5,0);
  % 尺寸：高 2
  \draw[<->, line width=0.45pt] (-0.55,0) -- (-0.55,3);
  \draw[line width=0.4pt] (-0.75,0) -- (-0.35,0);
  \draw[line width=0.4pt] (-0.75,3) -- (-0.35,3);
  \node[left] at (-0.60,1.5) {$2$};
  % 尺寸：底 2
  \draw[<->, line width=0.45pt] (0,-0.55) -- (3,-0.55);
  \draw[line width=0.4pt] (0,-0.75) -- (0,-0.35);
  \draw[line width=0.4pt] (3,-0.75) -- (3,-0.35);
  \node[below] at (1.5,-0.52) {$2$};
  \node at (1.5,-1.45) {正（主）视图};

  % ---------- 侧（左）视图 ----------
  \draw (6.4,0) -- (7.9,0) -- (6.4,3) -- cycle;
  \draw[<->, line width=0.45pt] (6.4,-0.55) -- (7.9,-0.55);
  \draw[line width=0.4pt] (6.4,-0.75) -- (6.4,-0.35);
  \draw[line width=0.4pt] (7.9,-0.75) -- (7.9,-0.35);
  \node[below] at (7.15,-0.52) {$1$};
  \node at (7.15,-1.45) {侧（左）视图};

  % ---------- 俯视图 ----------
  \draw (0,-3.05) -- (3,-3.05) -- (1.5,-4.55) -- cycle;
  % 尺寸：1 + 1
  \draw[<->, line width=0.45pt] (0,-2.55) -- (1.5,-2.55);
  \draw[<->, line width=0.45pt] (1.5,-2.55) -- (3,-2.55);
  \draw[line width=0.4pt] (0,-3.25) -- (0,-2.35);
  \draw[line width=0.4pt] (1.5,-3.25) -- (1.5,-2.35);
  \draw[line width=0.4pt] (3,-3.25) -- (3,-2.35);
  \node[above] at (0.75,-2.55) {$1$};
  \node[above] at (2.25,-2.55) {$1$};
  \node at (1.5,-5.05) {俯视图};
\end{tikzpicture}
